\documentclass[12pt,a4paper]{article}
\usepackage[utf8]{vietnam}
\usepackage{amsmath,amssymb}
\usepackage{tabularx}
\usepackage{tikz}
\usepackage{xcolor}
\usepackage[top=0.6cm, bottom=0.6cm, left=0.6cm, right=0.6cm]{geometry}
\usepackage[dethi]{ex_test}

\SAOPTN{widthSA=2.6,heightSA=0.62}
\setlength{\parskip}{0pt}

% Dong ke cham mo chuan viet tay
\newcommand{\dongke}[1]{%
  \begingroup%
  \setlength{\parskip}{0pt}%
  \baselineskip=18.5pt%
  \vspace{1.0mm}%
  \foreach \i in {1,...,#1}{%
    \noindent{\color{gray!60}\dotfill}\par%
  }%
  \vspace{0.6mm}%
  \endgroup%
}

% Dinh nghia o bai tap tu luan trong bang
\newcommand{\oBaiTap}[3][6]{%
  \vspace{1.0mm}%
  \noindent\textbf{#2} #3\par%
  \dongke{#1}%
}

\begin{document}
\pagestyle{empty}
% =============================================================================
% ========================= MA DE 101 - MAT TRUOC (TRANG 1/2) =============
% =============================================================================
\noindent
\begin{tabularx}{\linewidth}{|X|c|}
\hline
\textbf{PHIẾU KIỂM TRA ĐỊNH KỲ TOÁN 8 - TUẦN 07 - BUỔI 01} & \textbf{Mã đề: 101 - Trang 1/2} \\
\hline
\end{tabularx}

\vspace{1.2mm}

\noindent
\begin{tabularx}{\linewidth}{|p{3.2cm}|X|}
\hline
\centering\textbf{Điểm số} & \centering\textbf{Lời nhận xét của thầy cô} \tabularnewline
\rule{0pt}{1.8cm} & \\
\hline
\end{tabularx}

\vspace{1.2mm}

\noindent\textbf{Câu 1 (0,25 điểm).} Kết quả của $3x^2 y + 5x^2 y$ là: \shortans{$8x^2 y$}\par\vspace{0.3mm}
\noindent\textbf{Câu 2 (0,25 điểm).} Bậc của đa thức $A = 2x^3 y - 4xy + x^3 y + 5$ là: \shortans{$4$}\par\vspace{0.3mm}
\noindent\textbf{Câu 3 (0,25 điểm).} Giá trị của biểu thức $P = x^2 - 4xy$ tại $x = 2, y = -1$ là: \shortans{$12$}\par\vspace{0.3mm}
\noindent\textbf{Câu 4 (0,25 điểm).} Khai triển biểu thức $(x - 3)^2$ ta được $x^2 - \ldots + 9$. Điền vào chỗ trống: \shortans{$6x$}\par\vspace{0.3mm}
\noindent\textbf{Câu 5 (0,25 điểm).} Khai triển $x^2 - 25$ dưới dạng tích là: \shortans{$(x-5)(x+5)$}\par\vspace{0.3mm}
\noindent\textbf{Câu 6 (0,25 điểm).} Phân tích đa thức $5x^2 - 10x$ thành nhân tử được: \shortans{$5x(x-2)$}\par\vspace{0.3mm}
\noindent\textbf{Câu 7 (0,25 điểm).} Tứ giác $ABCD$ có $AB \parallel CD$ được gọi là hình gì? \shortans{hình thang}\par\vspace{0.3mm}
\noindent\textbf{Câu 8 (0,25 điểm).} Cho hình thang cân $ABCD$ ($AB \parallel CD$) có $\widehat{D} = 70^\circ$. Số đo góc $\widehat{C}$ là: \shortans{$70^\circ$}\par\vspace{0.3mm}
\noindent\textbf{Câu 9 (0,25 điểm).} Cho hình thang cân $ABCD$ có đường chéo $AC = 8\text{ cm}$. Độ dài đường chéo $BD$ là: \shortans{$8\text{ cm}$}\par\vspace{0.3mm}
\noindent\textbf{Câu 10 (0,25 điểm).} Một mảnh đất hình thang có hai đáy dài $10\text{ m}$ và $14\text{ m}$, chiều cao $5\text{ m}$. Diện tích mảnh đất là: \shortans{$60\text{ m}^2$}\par\vspace{0.3mm}

\vspace{0.5mm}
\noindent\textbf{Câu 11 (2,0 điểm). Thực hiện phép tính (hợp lý nếu có thể):}
\vspace{0.5mm}

\noindent
\begin{tabularx}{\linewidth}{|X|X|}
\hline
\oBaiTap[10]{a)}{$2x(3x^2 - x + 4) - 6x^3$} & \oBaiTap[10]{b)}{$(2x + 1)^2 - 4x(x - 2)$} \\
\hline
\oBaiTap[10]{c)}{$(15x^3 y^2 - 10x^2 y + 5xy) : 5xy$} & \oBaiTap[10]{d)}{$(x - 3)(x + 3) - (x + 2)(x - 5)$} \\
\hline
\end{tabularx}

\newpage
% =============================================================================
% ========================= MA DE 101 - MAT SAU (TRANG 2/2) ==============
% =============================================================================
\noindent
\begin{tabularx}{\linewidth}{|X|}
\hline
\centering\textbf{TOÁN 8 - TUẦN 07 - BUỔI 01 | Mã đề: 101 - Trang 2/2} \\
\hline
\end{tabularx}

\vspace{1.2mm}

\noindent\textbf{Câu 12 (2,0 điểm). Tìm $x$, biết:}
\vspace{0.5mm}

\noindent
\begin{tabularx}{\linewidth}{|X|X|}
\hline
\oBaiTap[5]{a)}{$3x(x - 2) - 3x^2 = 12$} & \oBaiTap[5]{b)}{$x(x - 3) + 4(x - 3) = 0$} \\
\hline
\oBaiTap[5]{c)}{$x^2 - 16 = 0$} & \oBaiTap[5]{d)}{$(x - 1)^2 - 4 = 0$} \\
\hline
\end{tabularx}

\vspace{1.2mm}

\noindent
\begin{tabularx}{\linewidth}{|X|}
\hline
\vspace{0.5mm}
\textbf{Câu 13 (2,5 điểm).} Cho hình thang cân $ABCD$ ($AB \parallel CD$, $AB < CD$). Kẻ các đường cao $AH \perp CD$ và $BK \perp CD$ ($H, K \in CD$).\par
\textbf{a) (1,25 điểm)} Chứng minh $\Delta AHD = \Delta BKC$ và suy ra $DH = CK$.\par
\textbf{b) (1,25 điểm)} Biết $AB = 6\text{ cm}$ và $CD = 14\text{ cm}$. Tính độ dài các đoạn thẳng $DH$ và $HC$.\par
\vspace{0.4mm}
\noindent\textit{Bài giải:}\dongke{11}
\\ \hline
\end{tabularx}

\vspace{1.2mm}

\noindent
\begin{tabularx}{\linewidth}{|X|}
\hline
\vspace{0.5mm}
\textbf{Câu 14 (1,0 điểm).} Cho hai số thực $a, b$ thỏa mãn $a + b = 5$ và $ab = 6$. Tính giá trị của biểu thức $P = a^3 + b^3$.\par
\vspace{0.4mm}
\noindent\textit{Bài giải:}\dongke{8}
\\ \hline
\end{tabularx}

\newpage
% =============================================================================
% ========================= MA DE 102 - MAT TRUOC (TRANG 1/2) =============
% =============================================================================
\noindent
\begin{tabularx}{\linewidth}{|X|c|}
\hline
\textbf{PHIẾU KIỂM TRA ĐỊNH KỲ TOÁN 8 - TUẦN 07 - BUỔI 01} & \textbf{Mã đề: 102 - Trang 1/2} \\
\hline
\end{tabularx}

\vspace{1.2mm}

\noindent
\begin{tabularx}{\linewidth}{|p{3.2cm}|X|}
\hline
\centering\textbf{Điểm số} & \centering\textbf{Lời nhận xét của thầy cô} \tabularnewline
\rule{0pt}{1.8cm} & \\
\hline
\end{tabularx}

\vspace{1.2mm}

\noindent\textbf{Câu 1 (0,25 điểm).} Kết quả của $4x^3 y + 6x^3 y$ là: \shortans{$10x^3 y$}\par\vspace{0.3mm}
\noindent\textbf{Câu 2 (0,25 điểm).} Bậc của đa thức $B = 3x^2 y^2 - 5x^3 + 2x^2 y^2 + 1$ là: \shortans{$4$}\par\vspace{0.3mm}
\noindent\textbf{Câu 3 (0,25 điểm).} Giá trị của biểu thức $Q = 2x^2 - 3xy$ tại $x = 1, y = -2$ là: \shortans{$8$}\par\vspace{0.3mm}
\noindent\textbf{Câu 4 (0,25 điểm).} Khai triển biểu thức $(x - 4)^2$ ta được $x^2 - \ldots + 16$. Điền vào chỗ trống: \shortans{$8x$}\par\vspace{0.3mm}
\noindent\textbf{Câu 5 (0,25 điểm).} Khai triển $x^2 - 36$ dưới dạng tích là: \shortans{$(x-6)(x+6)$}\par\vspace{0.3mm}
\noindent\textbf{Câu 6 (0,25 điểm).} Phân tích đa thức $6x^2 - 12x$ thành nhân tử được: \shortans{$6x(x-2)$}\par\vspace{0.3mm}
\noindent\textbf{Câu 7 (0,25 điểm).} Hình thang là tứ giác có hai cạnh đối như thế nào với nhau? \shortans{song song}\par\vspace{0.3mm}
\noindent\textbf{Câu 8 (0,25 điểm).} Cho hình thang cân $ABCD$ ($AB \parallel CD$) có $\widehat{D} = 65^\circ$. Số đo góc $\widehat{C}$ là: \shortans{$65^\circ$}\par\vspace{0.3mm}
\noindent\textbf{Câu 9 (0,25 điểm).} Cho hình thang cân $ABCD$ có đường chéo $AC = 10\text{ cm}$. Độ dài đường chéo $BD$ là: \shortans{$10\text{ cm}$}\par\vspace{0.3mm}
\noindent\textbf{Câu 10 (0,25 điểm).} Một mảnh đất hình thang có hai đáy dài $8\text{ m}$ và $12\text{ m}$, chiều cao $6\text{ m}$. Diện tích mảnh đất là: \shortans{$60\text{ m}^2$}\par\vspace{0.3mm}

\vspace{0.5mm}
\noindent\textbf{Câu 11 (2,0 điểm). Thực hiện phép tính (hợp lý nếu có thể):}
\vspace{0.5mm}

\noindent
\begin{tabularx}{\linewidth}{|X|X|}
\hline
\oBaiTap[10]{a)}{$3x(2x^2 - x + 3) - 6x^3$} & \oBaiTap[10]{b)}{$(2x + 3)^2 - 4x(x - 1)$} \\
\hline
\oBaiTap[10]{c)}{$(18x^3 y^2 - 12x^2 y + 6xy) : 6xy$} & \oBaiTap[10]{d)}{$(x - 4)(x + 4) - (x + 1)(x - 6)$} \\
\hline
\end{tabularx}

\newpage
% =============================================================================
% ========================= MA DE 102 - MAT SAU (TRANG 2/2) ==============
% =============================================================================
\noindent
\begin{tabularx}{\linewidth}{|X|}
\hline
\centering\textbf{TOÁN 8 - TUẦN 07 - BUỔI 01 | Mã đề: 102 - Trang 2/2} \\
\hline
\end{tabularx}

\vspace{1.2mm}

\noindent\textbf{Câu 12 (2,0 điểm). Tìm $x$, biết:}
\vspace{0.5mm}

\noindent
\begin{tabularx}{\linewidth}{|X|X|}
\hline
\oBaiTap[5]{a)}{$4x(x - 2) - 4x^2 = 16$} & \oBaiTap[5]{b)}{$x(x - 4) + 3(x - 4) = 0$} \\
\hline
\oBaiTap[5]{c)}{$x^2 - 25 = 0$} & \oBaiTap[5]{d)}{$(x - 2)^2 - 9 = 0$} \\
\hline
\end{tabularx}

\vspace{1.2mm}

\noindent
\begin{tabularx}{\linewidth}{|X|}
\hline
\vspace{0.5mm}
\textbf{Câu 13 (2,5 điểm).} Cho hình thang cân $ABCD$ ($AB \parallel CD$, $AB < CD$). Kẻ các đường cao $AH \perp CD$ và $BK \perp CD$ ($H, K \in CD$).\par
\textbf{a) (1,25 điểm)} Chứng minh $\Delta AHD = \Delta BKC$ và suy ra $DH = CK$.\par
\textbf{b) (1,25 điểm)} Biết $AB = 8\text{ cm}$ và $CD = 18\text{ cm}$. Tính độ dài các đoạn thẳng $DH$ và $HC$.\par
\vspace{0.4mm}
\noindent\textit{Bài giải:}\dongke{11}
\\ \hline
\end{tabularx}

\vspace{1.2mm}

\noindent
\begin{tabularx}{\linewidth}{|X|}
\hline
\vspace{0.5mm}
\textbf{Câu 14 (1,0 điểm).} Cho hai số thực $a, b$ thỏa mãn $a + b = 6$ và $ab = 8$. Tính giá trị của biểu thức $Q = a^3 + b^3$.\par
\vspace{0.4mm}
\noindent\textit{Bài giải:}\dongke{8}
\\ \hline
\end{tabularx}

\newpage
% =============================================================================
% ========================= MA DE 103 - MAT TRUOC (TRANG 1/2) =============
% =============================================================================
\noindent
\begin{tabularx}{\linewidth}{|X|c|}
\hline
\textbf{PHIẾU KIỂM TRA ĐỊNH KỲ TOÁN 8 - TUẦN 07 - BUỔI 01} & \textbf{Mã đề: 103 - Trang 1/2} \\
\hline
\end{tabularx}

\vspace{1.2mm}

\noindent
\begin{tabularx}{\linewidth}{|p{3.2cm}|X|}
\hline
\centering\textbf{Điểm số} & \centering\textbf{Lời nhận xét của thầy cô} \tabularnewline
\rule{0pt}{1.8cm} & \\
\hline
\end{tabularx}

\vspace{1.2mm}

\noindent\textbf{Câu 1 (0,25 điểm).} Kết quả của $2x y^2 + 7x y^2$ là: \shortans{$9x y^2$}\par\vspace{0.3mm}
\noindent\textbf{Câu 2 (0,25 điểm).} Bậc của đa thức $C = 4x y^3 - 2x^2 + x y^3 - 7$ là: \shortans{$4$}\par\vspace{0.3mm}
\noindent\textbf{Câu 3 (0,25 điểm).} Giá trị của biểu thức $R = 3x^2 - 2xy$ tại $x = 2, y = 1$ là: \shortans{$8$}\par\vspace{0.3mm}
\noindent\textbf{Câu 4 (0,25 điểm).} Khai triển biểu thức $(x - 5)^2$ ta được $x^2 - \ldots + 25$. Điền vào chỗ trống: \shortans{$10x$}\par\vspace{0.3mm}
\noindent\textbf{Câu 5 (0,25 điểm).} Khai triển $x^2 - 49$ dưới dạng tích là: \shortans{$(x-7)(x+7)$}\par\vspace{0.3mm}
\noindent\textbf{Câu 6 (0,25 điểm).} Phân tích đa thức $7x^2 - 14x$ thành nhân tử được: \shortans{$7x(x-2)$}\par\vspace{0.3mm}
\noindent\textbf{Câu 7 (0,25 điểm).} Tứ giác có hai cạnh đối song song gọi là: \shortans{hình thang}\par\vspace{0.3mm}
\noindent\textbf{Câu 8 (0,25 điểm).} Cho hình thang cân $ABCD$ ($AB \parallel CD$) có $\widehat{D} = 80^\circ$. Số đo góc $\widehat{C}$ là: \shortans{$80^\circ$}\par\vspace{0.3mm}
\noindent\textbf{Câu 9 (0,25 điểm).} Cho hình thang cân $ABCD$ có đường chéo $AC = 12\text{ cm}$. Độ dài đường chéo $BD$ là: \shortans{$12\text{ cm}$}\par\vspace{0.3mm}
\noindent\textbf{Câu 10 (0,25 điểm).} Một mảnh đất hình thang có hai đáy dài $6\text{ m}$ và $14\text{ m}$, chiều cao $7\text{ m}$. Diện tích mảnh đất là: \shortans{$70\text{ m}^2$}\par\vspace{0.3mm}

\vspace{0.5mm}
\noindent\textbf{Câu 11 (2,0 điểm). Thực hiện phép tính (hợp lý nếu có thể):}
\vspace{0.5mm}

\noindent
\begin{tabularx}{\linewidth}{|X|X|}
\hline
\oBaiTap[10]{a)}{$4x(2x^2 - x + 2) - 8x^3$} & \oBaiTap[10]{b)}{$(2x + 2)^2 - 4x(x - 3)$} \\
\hline
\oBaiTap[10]{c)}{$(21x^3 y^2 - 14x^2 y + 7xy) : 7xy$} & \oBaiTap[10]{d)}{$(x - 5)(x + 5) - (x + 3)(x - 4)$} \\
\hline
\end{tabularx}

\newpage
% =============================================================================
% ========================= MA DE 103 - MAT SAU (TRANG 2/2) ==============
% =============================================================================
\noindent
\begin{tabularx}{\linewidth}{|X|}
\hline
\centering\textbf{TOÁN 8 - TUẦN 07 - BUỔI 01 | Mã đề: 103 - Trang 2/2} \\
\hline
\end{tabularx}

\vspace{1.2mm}

\noindent\textbf{Câu 12 (2,0 điểm). Tìm $x$, biết:}
\vspace{0.5mm}

\noindent
\begin{tabularx}{\linewidth}{|X|X|}
\hline
\oBaiTap[5]{a)}{$5x(x - 2) - 5x^2 = 20$} & \oBaiTap[5]{b)}{$x(x - 5) + 2(x - 5) = 0$} \\
\hline
\oBaiTap[5]{c)}{$x^2 - 36 = 0$} & \oBaiTap[5]{d)}{$(x - 3)^2 - 16 = 0$} \\
\hline
\end{tabularx}

\vspace{1.2mm}

\noindent
\begin{tabularx}{\linewidth}{|X|}
\hline
\vspace{0.5mm}
\textbf{Câu 13 (2,5 điểm).} Cho hình thang cân $ABCD$ ($AB \parallel CD$, $AB < CD$). Kẻ các đường cao $AH \perp CD$ và $BK \perp CD$ ($H, K \in CD$).\par
\textbf{a) (1,25 điểm)} Chứng minh $\Delta AHD = \Delta BKC$ và suy ra $DH = CK$.\par
\textbf{b) (1,25 điểm)} Biết $AB = 4\text{ cm}$ và $CD = 16\text{ cm}$. Tính độ dài các đoạn thẳng $DH$ và $HC$.\par
\vspace{0.4mm}
\noindent\textit{Bài giải:}\dongke{11}
\\ \hline
\end{tabularx}

\vspace{1.2mm}

\noindent
\begin{tabularx}{\linewidth}{|X|}
\hline
\vspace{0.5mm}
\textbf{Câu 14 (1,0 điểm).} Cho hai số thực $a, b$ thỏa mãn $a + b = 7$ và $ab = 10$. Tính giá trị của biểu thức $R = a^3 + b^3$.\par
\vspace{0.4mm}
\noindent\textit{Bài giải:}\dongke{8}
\\ \hline
\end{tabularx}

\newpage
% =============================================================================
% ========================= MA DE 104 - MAT TRUOC (TRANG 1/2) =============
% =============================================================================
\noindent
\begin{tabularx}{\linewidth}{|X|c|}
\hline
\textbf{PHIẾU KIỂM TRA ĐỊNH KỲ TOÁN 8 - TUẦN 07 - BUỔI 01} & \textbf{Mã đề: 104 - Trang 1/2} \\
\hline
\end{tabularx}

\vspace{1.2mm}

\noindent
\begin{tabularx}{\linewidth}{|p{3.2cm}|X|}
\hline
\centering\textbf{Điểm số} & \centering\textbf{Lời nhận xét của thầy cô} \tabularnewline
\rule{0pt}{1.8cm} & \\
\hline
\end{tabularx}

\vspace{1.2mm}

\noindent\textbf{Câu 1 (0,25 điểm).} Kết quả của $5x^2 y^2 + 3x^2 y^2$ là: \shortans{$8x^2 y^2$}\par\vspace{0.3mm}
\noindent\textbf{Câu 2 (0,25 điểm).} Bậc của đa thức $D = 5x^4 y - 3x^3 + 2x^4 y + 6$ là: \shortans{$5$}\par\vspace{0.3mm}
\noindent\textbf{Câu 3 (0,25 điểm).} Giá trị của biểu thức $S = x^2 - 5xy$ tại $x = -1, y = 2$ là: \shortans{$11$}\par\vspace{0.3mm}
\noindent\textbf{Câu 4 (0,25 điểm).} Khai triển biểu thức $(x - 6)^2$ ta được $x^2 - \ldots + 36$. Điền vào chỗ trống: \shortans{$12x$}\par\vspace{0.3mm}
\noindent\textbf{Câu 5 (0,25 điểm).} Khai triển $x^2 - 64$ dưới dạng tích là: \shortans{$(x-8)(x+8)$}\par\vspace{0.3mm}
\noindent\textbf{Câu 6 (0,25 điểm).} Phân tích đa thức $8x^2 - 16x$ thành nhân tử được: \shortans{$8x(x-2)$}\par\vspace{0.3mm}
\noindent\textbf{Câu 7 (0,25 điểm).} Trong hình thang $MNPQ$ ($MN \parallel PQ$), hai đoạn thẳng $MN$ và $PQ$ gọi là: \shortans{hai đáy}\par\vspace{0.3mm}
\noindent\textbf{Câu 8 (0,25 điểm).} Cho hình thang cân $ABCD$ ($AB \parallel CD$) có $\widehat{D} = 75^\circ$. Số đo góc $\widehat{C}$ là: \shortans{$75^\circ$}\par\vspace{0.3mm}
\noindent\textbf{Câu 9 (0,25 điểm).} Cho hình thang cân $ABCD$ có đường chéo $AC = 15\text{ cm}$. Độ dài đường chéo $BD$ là: \shortans{$15\text{ cm}$}\par\vspace{0.3mm}
\noindent\textbf{Câu 10 (0,25 điểm).} Một mảnh đất hình thang có hai đáy dài $12\text{ m}$ và $16\text{ m}$, chiều cao $5\text{ m}$. Diện tích mảnh đất là: \shortans{$70\text{ m}^2$}\par\vspace{0.3mm}

\vspace{0.5mm}
\noindent\textbf{Câu 11 (2,0 điểm). Thực hiện phép tính (hợp lý nếu có thể):}
\vspace{0.5mm}

\noindent
\begin{tabularx}{\linewidth}{|X|X|}
\hline
\oBaiTap[10]{a)}{$2x(4x^2 - 2x + 3) - 8x^3$} & \oBaiTap[10]{b)}{$(2x + 4)^2 - 4x(x - 2)$} \\
\hline
\oBaiTap[10]{c)}{$(24x^3 y^2 - 16x^2 y + 8xy) : 8xy$} & \oBaiTap[10]{d)}{$(x - 2)(x + 2) - (x + 4)(x - 3)$} \\
\hline
\end{tabularx}

\newpage
% =============================================================================
% ========================= MA DE 104 - MAT SAU (TRANG 2/2) ==============
% =============================================================================
\noindent
\begin{tabularx}{\linewidth}{|X|}
\hline
\centering\textbf{TOÁN 8 - TUẦN 07 - BUỔI 01 | Mã đề: 104 - Trang 2/2} \\
\hline
\end{tabularx}

\vspace{1.2mm}

\noindent\textbf{Câu 12 (2,0 điểm). Tìm $x$, biết:}
\vspace{0.5mm}

\noindent
\begin{tabularx}{\linewidth}{|X|X|}
\hline
\oBaiTap[5]{a)}{$2x(x - 3) - 2x^2 = 12$} & \oBaiTap[5]{b)}{$x(x - 6) + 5(x - 6) = 0$} \\
\hline
\oBaiTap[5]{c)}{$x^2 - 49 = 0$} & \oBaiTap[5]{d)}{$(x - 4)^2 - 25 = 0$} \\
\hline
\end{tabularx}

\vspace{1.2mm}

\noindent
\begin{tabularx}{\linewidth}{|X|}
\hline
\vspace{0.5mm}
\textbf{Câu 13 (2,5 điểm).} Cho hình thang cân $ABCD$ ($AB \parallel CD$, $AB < CD$). Kẻ các đường cao $AH \perp CD$ và $BK \perp CD$ ($H, K \in CD$).\par
\textbf{a) (1,25 điểm)} Chứng minh $\Delta AHD = \Delta BKC$ và suy ra $DH = CK$.\par
\textbf{b) (1,25 điểm)} Biết $AB = 10\text{ cm}$ và $CD = 22\text{ cm}$. Tính độ dài các đoạn thẳng $DH$ và $HC$.\par
\vspace{0.4mm}
\noindent\textit{Bài giải:}\dongke{11}
\\ \hline
\end{tabularx}

\vspace{1.2mm}

\noindent
\begin{tabularx}{\linewidth}{|X|}
\hline
\vspace{0.5mm}
\textbf{Câu 14 (1,0 điểm).} Cho hai số thực $a, b$ thỏa mãn $a + b = 4$ và $ab = 3$. Tính giá trị của biểu thức $S = a^3 + b^3$.\par
\vspace{0.4mm}
\noindent\textit{Bài giải:}\dongke{8}
\\ \hline
\end{tabularx}

\newpage
\end{document}
